\documentclass{article}
\usepackage{mathtools} 
\usepackage{fontspec}
\usepackage[UTF8]{ctex}
\usepackage{amsthm}
\usepackage{mdframed}
\usepackage{xcolor}
\usepackage{amssymb}
\usepackage{amsmath}


% 定义新的带灰色背景的说明环境 zremark
\newmdtheoremenv[
  backgroundcolor=gray!10,
  % 边框与背景一致，边框线会消失
  linecolor=gray!10
]{zremark}{说明}

% 通用矩阵命令: \flexmatrix{矩阵名}{元素符号}{行数}{列数}
\newcommand{\flexmatrix}[4]{
  \[
  #1 = \begin{pmatrix}
    #2_{11}     & #2_{12}     & \cdots & #2_{1#4}   \\
    #2_{21}     & #2_{22}     & \cdots & #2_{2#4}   \\
    \vdots      & \vdots      & \ddots & \vdots     \\
    #2_{#31}    & #2_{#32}    & \cdots & #2_{#3#4}
  \end{pmatrix}
  \]
}

% 简化版命令（默认矩阵名为A，元素符号为a）: \quickmatrix{行数}{列数}
\newcommand{\quickmatrix}[2]{\flexmatrix{A}{a}{#1}{#2}}

\begin{document}
\title{注释 8.2}
\author{张志聪}
\maketitle

\begin{zremark}
  % （第二换元积分法）设函数$f(x)$在区间$I$上有定义，$\varphi(t)$在区间$J$上可导，
  % $\varphi(J) = I$，且$x = \varphi(t)$在区间$J$上存在反函数$t = \varphi^{-1}(x), x \in I$。
  % 如果不定积分$\int f(x) dx$在$I$上存在，则当不定积分$\int f(\varphi(t)) \varphi'(t) dt = G(t) + C$
  % 在$J$上存在时，在$I$上有
  % \begin{align*}
  %   \int f(x) dx = G(\varphi^{-1}(x)) + C
  % \end{align*}
  为什么说“$\varphi'(t) \neq 0, x \in J$”是比“$x = \varphi(t)$在区间$J$上存在反函数$t = \varphi^{-1}(x), x \in I$”
  更强的条件。

  书中有些步骤省略了，这里做些补充。
\end{zremark}

先纠正书中的错误点：
“$\varphi'(t) \neq 0, x \in J$”，应该是“$\varphi'(t) \neq 0, x \in I$”。

\textbf{证明：}

因为“$\varphi'(t) \neq 0, x \in I$”，
$\varphi'(t)$在$J$上是同号的。反证法，假设存在$t_1, t_2 \in J$，使得
\begin{align*}
  \varphi'(t_1) > 0 \\
  \varphi'(t_2) < 0
\end{align*}
由定理6.5（达布（Darboux）定理）可知，存在$\xi \in (t_1, t_2)$，使得
\begin{align*}
  \varphi'(\xi) = 0
\end{align*}
这与题设矛盾。

所以$\varphi(t)$在$J$上是严格单调的（定理6.4 推论），于是可得$\varphi(t)$在$J$上是可逆的，
且由于$\varphi(t) \neq 0$，所以反函数$t = \varphi^{-1}(x)$在$I$上可导，且
\begin{align*}
  (\varphi^{-1}(x))' = \frac{1}{\varphi'(t)}
\end{align*}

\begin{zremark}
  分部积分“绕回原式”现象。
\end{zremark}

若对 $\int u \, dv$ 使用一次分部积分，再对所得的 $\int v \, du$ 反向再分部（交换 $u$ 与 $v$），则：
\[
  \begin{aligned}
    I          & = \int u \, dv = uv - \int v \, du, \\
    \int v \, du & = uv - \int u \, dv = uv - I.
  \end{aligned}
\]
代回第一式：
\[
  I = uv - (uv - I) = I.
\]


\end{document}